(1) 动能由转台转动与人体质点转动组成：\(T=I\dot\theta^2/2+mR^2(\dot\theta+\dot\phi)^2/2\)。\(\theta\) 是循环坐标；惯性系总角动量 \(L_z=(I+mR^2)\dot\theta+mR^2\dot\phi\) 守恒。

(2) 初始 \(L_z=0\)，所以 \(\dot\theta=-[mR^2/(I+mR^2)]\Omega\)。人相对地面的角速度为 \(\dot\theta+\Omega=I\Omega/(I+mR^2)\)。绕平台相对一周 \(\Delta\phi=2\pi\)，转台转角 \(\Delta\theta=-2\pi mR^2/(I+mR^2)\)。

(3) 代入角动量约束，\(T=\frac12[I mR^2/(I+mR^2)]\Omega^2\)。由始静止的动能增加可知内部驱动净机械功至少为该值。\(I\gg mR^2\) 时人近乎以 \(\Omega\) 对地转动，平台几乎静止；\(I\ll mR^2\) 时转台近乎以 \(-\Omega\) 转动，人对地近乎静止。

(4) 最后若相对角速度 \(\dot\phi=0\)，角动量式变为 \((I+mR^2)\dot\theta=0\)，故两者均停止，理想纯机械动能为零。理想可逆制动可把此前输入的动能回收；现实中的脚底摩擦及肌肉代谢会把部分能量变为热，因此角动量守恒不等于机械能可回收。
